\documentclass[11pt,letterpaper]{article}

\usepackage[spanish,es-noshorthands]{babel}
\usepackage{fontspec}
\setmainfont{Source Serif Pro}
\usepackage{microtype}
\usepackage{geometry}
\usepackage{booktabs}
\usepackage{tabularx}
\usepackage{enumitem}
\usepackage{xcolor}
\usepackage{graphicx}
\usepackage{csquotes}
\usepackage{amsmath}
\usepackage{siunitx}
\usepackage{tikz}
\usetikzlibrary{arrows.meta,positioning,shapes.geometric,fit,backgrounds,patterns}
\usepackage{xurl}
\usepackage[hidelinks]{hyperref}
\usepackage[backend=biber,style=apa,sorting=nyt]{biblatex}
\DeclareLanguageMapping{spanish}{spanish-apa}

\geometry{margin=2.3cm}
% Los rótulos de las figuras usan nodos de ancho fijo con texto alineado a la
% derecha, y la bibliografía compone en modo sloppy: ambos generan avisos de
% línea holgada que son cosméticos y no indican un problema de composición.
% Se silencian para que `make check` siga siendo sensible a los desbordes reales.
\hbadness=10000
\setlength{\emergencystretch}{1em}
\setlength{\parindent}{0pt}
\setlength{\parskip}{0.55em}
\setlist{nosep}
\addbibresource{referencias.bib}
\AtBeginBibliography{\sloppy}

\sisetup{
  output-decimal-marker = {,},
  group-separator = {.},
  group-minimum-digits = 4,
  detect-all
}

% Colores: vehículo privado, transporte público, modos a pie, mi propio caso
\definecolor{cAuto}{RGB}{191,110,30}
\definecolor{cTP}{RGB}{31,78,121}
\definecolor{cPie}{RGB}{40,124,70}
\definecolor{cCR}{RGB}{150,40,55}
\definecolor{cGris}{RGB}{140,140,140}

% ---------------------------------------------------------------------------
% Barra horizontal de cuota modal. #1 fila, #2 longitud, #3 color, #4 rótulo,
% #5 valor. La escala es de 0,13 cm por punto porcentual.
% ---------------------------------------------------------------------------
\newcommand{\cuotabar}[5]{%
  \node[left, text width=4.6cm, align=right, font=\scriptsize] at (-0.18,#1) {#4};%
  \fill[#3] (0,#1-0.19) rectangle (#2,#1+0.19);%
  \node[right, font=\scriptsize] at (#2+0.12,#1) {#5};%
}

\title{\vspace{-1.2cm}\textbf{¿Un mundo sin automóviles? La respuesta no la
da el motor, la da el mapa}\\
\large Por qué la misma tecnología produce cuotas del 20 al 71 por ciento
según dónde se mida}
\author{Carlos Luis Rojas Aragonés\\
\small Gestión del Desarrollo Tecnológico}
\date{\small 19 de agosto de 2026}

\begin{document}

\maketitle
\vspace{-0.6em}

\section{Mi respuesta, y por qué no habla de mí}

Sí, conduzco. Y no, hoy no puedo imaginarme una vida sin vehículo privado en
la Gran Área Metropolitana de Costa Rica, donde vivo. Pero esa respuesta no
dice nada interesante sobre mis preferencias: dice cuánto mide la distancia
entre las cosas que necesito. En Costa Rica el 41,6 por ciento de los viajes
se hace en vehículo privado motorizado y el 49,8 en autobús
\parencite{allen2024}; en la almendra central de Madrid el vehículo privado
baja al 20,3 por ciento y en la corona regional de la misma comunidad autónoma
sube al 56,2 \parencite{crtm2019}. Mismo país, misma flota, mismo precio del
combustible, misma tecnología: la cuota se multiplica por 2,8 al cambiar de
código postal.

El enunciado ya contiene la observación clave, que su entorno influye e
incluso determina la respuesta, y conviene tomarla en serio hasta el final:
\textbf{si el mismo artefacto produce resultados que van del 20 al 71 por
ciento según dónde se mida, entonces la variable explicativa no es el
artefacto}. La pregunta \enquote{¿puede vivir sin coche?} está mal planteada
como pregunta tecnológica. Bien planteada es una pregunta de accesibilidad:
¿cuántos de mis destinos habituales caben dentro de mi presupuesto diario de
tiempo de viaje sin usar un coche?

Esa reformulación tiene una consecuencia incómoda para quien decide sobre
tecnología. Mejorar el automóvil, hacerlo eléctrico, autónomo o compartido, no
mueve la variable que decide. La cambia el uso del suelo, y el uso del suelo se
fija con decisiones de infraestructura cuyo plazo de maduración es de décadas.
En mi caso esa fecha ya está puesta en un calendario: la ley de financiamiento
del tren eléctrico de pasajeros de la GAM se firmó el 25 de mayo de 2026
\parencite{incofer2026} y el sistema completo no operará antes de 2031
\parencite{teletica2026}. Mi respuesta de hoy es \enquote{no puedo}, y la
fecha más temprana en que podría cambiar no la decido yo.

\section{Marco de análisis: la figura de mérito equivocada}

Una \textbf{figura de mérito} (FOM) es una medida cuantitativa del desempeño de
una tecnología en la función que presta \parencite{deweck2022}. El automóvil
gana en casi todas las FOM del artefacto: velocidad punta, autonomía,
protección frente a la lluvia, disponibilidad a cualquier hora, capacidad de
carga. Si la comparación se hiciera ahí, el debate estaría cerrado. La
discusión de este foro sólo tiene sentido porque la FOM que decide es otra, y
es del sistema, no del vehículo. Tres resultados de la literatura la delimitan.

\begin{enumerate}
  \item \textbf{El presupuesto de tiempo de viaje es aproximadamente
    constante.} \textcite{zahavi1974} documentó que los hogares asignan al
    desplazamiento una fracción estable de su tiempo, y \textcite{marchetti1994}
    la resumió en la cifra que hoy lleva su nombre: alrededor de una hora
    diaria por persona, un valor que reaparece en sociedades muy distintas.
    Si el tiempo es fijo y la velocidad sube, lo que crece no es el tiempo
    libre: es el radio. Y el área accesible crece con el cuadrado de la
    velocidad.
  \item \textbf{La capacidad vial genera su propia demanda.}
    \textcite{duranton2011} estimaron con datos de ciudades estadounidenses que
    los kilómetros recorridos crecen de forma aproximadamente proporcional a
    los kilómetros de carril construidos, con una elasticidad cercana a uno.
    Su conclusión es explícita: ampliar la red, e incluso ampliar el transporte
    público, no reduce la congestión. Este es el resultado que convierte
    \enquote{construir más carriles} en una estrategia sin punto de llegada.
  \item \textbf{Lo que explica cuánto se conduce es el entorno construido, no
    la preferencia.} El metaanálisis de \textcite{ewing2010} sobre decenas de
    estudios concluye que los kilómetros recorridos por habitante dependen
    sobre todo de la accesibilidad a destinos y, en segundo lugar, del diseño
    de la red de calles, y que caminar depende de la mezcla de usos y de la
    densidad de intersecciones. Es la formalización empírica de las tres
    \enquote{d} de \textcite{cervero1997}: densidad, diversidad y diseño.
\end{enumerate}

Juntando las tres piezas aparece el mecanismo completo, y es un bucle:
el automóvil multiplica el radio accesible en un tiempo dado; la ciudad usa ese
radio para separar la vivienda del trabajo y del comercio; una vez separados,
el automóvil deja de ser una opción y pasa a ser un requisito de acceso al
empleo. Es \textbf{dependencia de trayectoria} en el sentido estricto de
\textcite{david1985} y \textcite{arthur1989}: cada decisión individual
racional refuerza la configuración que hace obligatoria la siguiente. Lo que
\textcite{newman1999} llamaron dependencia del automóvil no es un hábito
cultural, es una propiedad geométrica del territorio ya construido.

\section{La evidencia: una escalera del 20 al 71 por ciento}

\begin{figure}[ht]
  \centering
  \begin{tikzpicture}
    % escala: 0,13 cm por punto porcentual
    \foreach \xv/\lab in {0/0, 2.6/20, 5.2/40, 7.8/60, 10.4/80} {
      \draw[gray!22] (\xv,-0.10) -- (\xv,5.85);
      \node[below, font=\tiny, gray!55!black] at (\xv,-0.12) {\lab\,\%};
    }
    \draw[red!65!black, thick, densely dotted] (6.5,-0.10) -- (6.5,5.85);
    \node[font=\tiny, red!65!black, above right] at (6.52,5.60)
      {la mitad de los desplazamientos};
    \cuotabar{5.40}{9.256}{cTP}{Los Ángeles, ciudad\\\tiny (al trabajo, ACS
      2024)}{71,2\,\%}
    \cuotabar{4.80}{8.775}{cTP}{Gran Adelaida\\\tiny (al trabajo, censo
      2021)}{67,5\,\%}
    \cuotabar{4.20}{7.306}{cAuto}{Madrid, corona regional\\\tiny (todos los
      viajes, EDM 2018)}{56,2\,\%}
    \cuotabar{3.60}{6.201}{cAuto}{Madrid, corona metropolitana\\\tiny (todos
      los viajes, EDM 2018)}{47,7\,\%}
    \cuotabar{3.00}{5.408}{cCR}{\textbf{Costa Rica}\\\tiny (todos los viajes,
      CGR 2023)}{\textbf{41,6\,\%}}
    \cuotabar{2.40}{5.070}{cAuto}{Comunidad de Madrid\\\tiny (todos los viajes,
      EDM 2018)}{39,0\,\%}
    \cuotabar{1.80}{4.212}{cAuto}{Madrid, periferia del municipio\\\tiny (todos
      los viajes, EDM 2018)}{32,4\,\%}
    \cuotabar{1.20}{3.263}{cTP}{Nueva York, ciudad\\\tiny (al trabajo, ACS
      2024)}{25,1\,\%}
    \cuotabar{0.60}{2.639}{cAuto}{Madrid, almendra central\\\tiny (todos los
      viajes, EDM 2018)}{20,3\,\%}
    \begin{scope}[shift={(0,-0.75)}]
      \fill[cAuto] (0,-0.09) rectangle (0.45,0.09);
      \node[right, font=\tiny] at (0.5,0) {cuota sobre todos los viajes};
      \fill[cTP] (4.30,-0.09) rectangle (4.75,0.09);
      \node[right, font=\tiny] at (4.80,0)
        {cuota sobre los desplazamientos al trabajo};
      \fill[cCR] (9.30,-0.09) rectangle (9.75,0.09);
      \node[right, font=\tiny] at (9.80,0) {mi entorno};
    \end{scope}
  \end{tikzpicture}
  \caption{Cuota del vehículo privado motorizado en nueve ámbitos. Las cuatro
  filas de Madrid y la de la Comunidad proceden de la encuesta domiciliaria de
  movilidad de 2018, sobre \num{15800000} desplazamientos diarios
  \parencite{crtm2019}; la de Costa Rica, de la Encuesta Nacional de Percepción
  de los Servicios Públicos de la Contraloría General de la República de 2023,
  recogida en \textcite{allen2024}; las de Nueva York y Los Ángeles, de la
  tabla B08301 de la American Community Survey de 2024 \parencite{acs2024};
  la de Gran Adelaida, del censo australiano de 2021, sumando conductor
  (63,3\,\%) y acompañante (4,2\,\%) \parencite{abs2021}. Las dos métricas no
  son intercambiables y por eso se distinguen con color; la discusión de esa
  limitación está en la sección de precisiones. Elaboración propia.}
  \label{fig:escalera}
\end{figure}

La Figura~\ref{fig:escalera} ordena las cuotas del vehículo privado en los
cuatro entornos que menciona el enunciado y en los míos. El resultado es una
escalera continua, sin escalón entre \enquote{ciudades donde se puede} y
\enquote{ciudades donde no}, y con un rango de 3,5 veces entre el extremo
inferior y el superior.

Dos lecturas conviene fijar antes de seguir.

\textbf{La primera es que Nueva York y Los Ángeles no se diferencian por la
tecnología disponible.} Son el mismo país, el mismo mercado de automóviles, el
mismo precio del combustible y prácticamente la misma regulación federal. En
Nueva York el 48,7 por ciento de los ocupados va al trabajo en transporte
público y el 25,1 en coche; en Los Ángeles esas cifras son 6,2 y 71,2
\parencite{acs2024}. La misma fuente da el dato patrimonial, que es todavía
más elocuente: el 56,7 por ciento de los hogares neoyorquinos no dispone de
ningún vehículo, frente al 12,0 por ciento de los angelinos. La diferencia
entre ambas ciudades es de forma urbana, no de artefacto.

\textbf{La segunda es que el enunciado acierta al no oponer ciudad y campo.}
Gran Adelaida es una capital estatal con 1,4 millones de habitantes y una cuota
del vehículo privado de 67,5 por ciento, superior a la de la corona regional de
Madrid, que es el ámbito más periférico y menos urbano de la muestra española.
Lo que separa a Adelaida de Madrid no es ser más o menos urbana, sino la
densidad y la mezcla de usos con que se construyó.

\section{El experimento natural de Madrid}

\begin{figure}[ht]
  \centering
  \begin{tikzpicture}
    % ---------------- panel A: reparto modal por corona ----------------
    \node[font=\scriptsize\itshape, gray!35!black, right] at (-4.30,3.78)
      {Panel A. Reparto modal según corona de residencia, día laborable};
    \foreach \xv/\lab in {0/0, 2.5/25, 5.0/50, 7.5/75, 10.0/100} {
      \draw[gray!22] (\xv,-0.15) -- (\xv,2.95);
      \node[below, font=\tiny, gray!55!black] at (\xv,-0.17) {\lab\,\%};
    }
    \foreach \y/\lab/\a/\b/\c in {%
      2.60/{Almendra central}/4.00/7.48/9.51,
      1.95/{Periferia del municipio}/3.22/6.50/9.74,
      1.30/{Corona metropolitana}/3.40/5.04/9.81,
      0.65/{Corona regional}/2.96/4.04/9.66,
      0.00/{\textbf{Comunidad de Madrid}}/3.40/5.83/9.73} {
      \node[left, font=\scriptsize] at (-0.12,\y) {\lab};
      \fill[cPie] (0,\y-0.20) rectangle (\a,\y+0.20);
      \fill[cTP] (\a,\y-0.20) rectangle (\b,\y+0.20);
      \fill[cAuto] (\b,\y-0.20) rectangle (\c,\y+0.20);
      \fill[gray!45] (\c,\y-0.20) rectangle (10.0,\y+0.20);
    }
    \node[font=\tiny, white] at (2.00,2.60) {40,0};
    \node[font=\tiny, white] at (5.74,2.60) {34,8};
    \node[font=\tiny, white] at (8.50,2.60) {20,3};
    \node[font=\tiny, white] at (1.61,1.95) {32,2};
    \node[font=\tiny, white] at (4.86,1.95) {32,8};
    \node[font=\tiny, white] at (8.12,1.95) {32,4};
    \node[font=\tiny, white] at (1.70,1.30) {34,0};
    \node[font=\tiny, white] at (4.22,1.30) {16,4};
    \node[font=\tiny, white] at (7.43,1.30) {47,7};
    \node[font=\tiny, white] at (1.48,0.65) {29,6};
    \node[font=\tiny, white] at (3.50,0.65) {10,8};
    \node[font=\tiny, white] at (6.85,0.65) {56,2};
    \node[font=\tiny, white] at (1.70,0.00) {34,0};
    \node[font=\tiny, white] at (4.62,0.00) {24,3};
    \node[font=\tiny, white] at (7.78,0.00) {39,0};
    \begin{scope}[shift={(0.10,3.25)}]
      \fill[cPie] (0,-0.09) rectangle (0.40,0.09);
      \node[right, font=\tiny] at (0.45,0) {a pie};
      \fill[cTP] (1.55,-0.09) rectangle (1.95,0.09);
      \node[right, font=\tiny] at (2.00,0) {transporte público};
      \fill[cAuto] (4.60,-0.09) rectangle (5.00,0.09);
      \node[right, font=\tiny] at (5.05,0) {vehículo privado};
      \fill[gray!45] (7.35,-0.09) rectangle (7.75,0.09);
      \node[right, font=\tiny] at (7.80,0) {otros modos};
    \end{scope}
    % ---------------- panel B: hogares sin turismo ----------------
    \begin{scope}[shift={(0,-4.15)}]
      \node[font=\scriptsize\itshape, gray!35!black, right] at (-4.30,3.30)
        {Panel B. Hogares que no tienen ningún turismo, misma encuesta};
      \foreach \xv/\lab in {0/0, 2.2/10, 4.4/20, 6.6/30, 8.8/40} {
        \draw[gray!22] (\xv,-0.15) -- (\xv,2.95);
        \node[below, font=\tiny, gray!55!black] at (\xv,-0.17) {\lab\,\%};
      }
      \foreach \y/\lab/\x/\v in {%
        2.60/{Almendra central}/8.272/{37,6},
        1.95/{Periferia del municipio}/6.600/{30,0},
        1.30/{Corona metropolitana}/3.718/{16,9},
        0.65/{Corona regional}/2.640/{12,0},
        0.00/{\textbf{Comunidad de Madrid}}/5.456/{24,8}} {
        \node[left, font=\scriptsize] at (-0.12,\y) {\lab};
        \fill[cTP!80] (0,\y-0.20) rectangle (\x,\y+0.20);
        \node[right, font=\scriptsize] at (\x+0.12,\y) {\v\,\%};
      }
      \draw[-{Stealth[length=2mm]}, thick, red!65!black]
        (9.90,2.60) -- (9.90,0.65);
      \node[right, font=\tiny, red!65!black, align=left, text width=2.4cm]
        at (9.98,1.62) {Hacia fuera, los hogares sin coche caen a un\\ tercio};
    \end{scope}
  \end{tikzpicture}
  \caption{El mismo territorio administrativo, la misma flota y el mismo
  precio del combustible, con cuatro resultados distintos. Datos de la encuesta
  domiciliaria de movilidad de la Comunidad de Madrid de 2018, publicada en
  noviembre de 2019 sobre \num{85064} encuestas y \num{15800000} viajes diarios
  \parencite{crtm2019}. La almendra central es el interior de la M-30. Los
  porcentajes del Panel A suman cien por fila; los de \enquote{otros modos}
  incluyen moto, bicicleta, taxi y transporte discrecional, que en el conjunto
  de la región representan el 2,7 por ciento. Elaboración propia.}
  \label{fig:madrid}
\end{figure}

La Figura~\ref{fig:madrid} es, a mi juicio, la pieza que resuelve el debate,
porque se acerca todo lo que permite el mundo real a un experimento
controlado. Dentro de una sola comunidad autónoma, con la misma oferta de
vehículos, los mismos impuestos, el mismo clima y la misma cultura, la cuota
del vehículo privado pasa del 20,3 al 56,2 por ciento y la proporción de
hogares sin ningún turismo cae del 37,6 al 12,0 por ciento
\parencite{crtm2019}. Lo único que cambia es la distancia a la que están las
cosas.

Dos detalles refuerzan la lectura. El primero es que en la almendra central el
transporte público (34,8 por ciento) supera al vehículo privado (20,3), y el
modo más usado es simplemente caminar (40,0): en un entorno denso y mezclado,
la alternativa dominante al coche no es otro vehículo, es la proximidad. El
segundo es que el 37,6 por ciento de los hogares del centro de Madrid ya vive
sin automóvil. No es una hipótesis de futuro ni una preferencia militante: es
la conducta observada de más de un tercio de los hogares cuando el entorno lo
permite.

\section{Por qué mejorar el automóvil no cambia la respuesta}

\begin{figure}[ht]
  \centering
  \begin{tikzpicture}
    % ---------------- panel A: el radio de media hora ----------------
    \node[font=\scriptsize\itshape, gray!35!black, right] at (-2.20,6.45)
      {Panel A. Lo que compran treinta minutos de viaje a cada velocidad};
    \begin{scope}[shift={(-2.20,0)}]
      \draw[gray!55!black] (0,0) -- (6.60,0);
      \draw[gray!55!black] (0,0) -- (0,6.10);
      \foreach \r/\col in {0.5/cPie, 1.5/cPie, 3.0/cTP, 6.0/cAuto} {
        \draw[\col, thick] (\r,0) arc (0:90:\r);
      }
      \fill[cAuto!10] (0,0) -- (6.0,0) arc (0:90:6.0) -- cycle;
      \fill[cTP!14] (0,0) -- (3.0,0) arc (0:90:3.0) -- cycle;
      \fill[cPie!22] (0,0) -- (1.5,0) arc (0:90:1.5) -- cycle;
      \fill[cPie!45] (0,0) -- (0.5,0) arc (0:90:0.5) -- cycle;
      \foreach \r/\col in {0.5/cPie, 1.5/cPie, 3.0/cTP, 6.0/cAuto} {
        \draw[\col, thick] (\r,0) arc (0:90:\r);
      }
      \node[font=\tiny, cPie, right] at (0.40,0.42)
        {a pie, 5 km/h: 2,5 km};
      \node[font=\tiny, cPie, right] at (1.14,1.18)
        {bicicleta, 15 km/h: 7,5 km};
      \node[font=\tiny, cTP, right] at (2.24,2.28)
        {autobús, 30 km/h: 15 km};
      \node[font=\tiny, cAuto, right] at (4.37,4.40)
        {automóvil, 60 km/h: 30 km};
      \node[font=\scriptsize, align=left, right] at (6.75,5.10)
        {A presupuesto de tiempo constante,\\ el radio crece con la velocidad
         y\\ el área con su cuadrado:};
      \node[font=\scriptsize, align=left, right] at (6.75,3.90)
        {$A = \pi (v\,t)^2$, con $t = 0{,}5$ h};
      \node[font=\scriptsize, align=left, right] at (6.75,2.95)
        {a pie: \SI{20}{\kilo\metre\squared}\\
         bicicleta: \SI{177}{\kilo\metre\squared}\\
         autobús: \SI{707}{\kilo\metre\squared}\\
         automóvil: \SI{2827}{\kilo\metre\squared}};
      \node[font=\scriptsize, align=left, right, red!65!black] at (6.75,1.30)
        {El automóvil multiplicó por 144 el\\ área alcanzable frente a caminar.
         \\ La ciudad gastó ese factor en\\ separarse, no en ahorrar tiempo.};
    \end{scope}
    % ---------------- panel B: capacidad de un carril ----------------
    \begin{scope}[shift={(0,-3.30)}]
      \node[font=\scriptsize\itshape, gray!35!black, right] at (0,2.55)
        {Panel B. Personas por hora que puede mover un carril, según a qué se
         dedique};
      \foreach \xv/\lab in {0/0, 1.8/{5.000}, 3.6/{10.000}, 5.4/{15.000},
                            7.2/{20.000}, 9.0/{25.000}} {
        \draw[gray!22] (\xv,-0.35) -- (\xv,2.20);
        \node[below, font=\tiny, gray!55!black] at (\xv,-0.37) {\lab};
      }
      \node[left, text width=3.5cm, align=right, font=\scriptsize]
        at (-0.15,1.90) {Tráfico general de vehículos privados};
      \fill[cAuto] (0,1.71) rectangle (0.216,2.09);
      \fill[cAuto!45] (0.216,1.71) rectangle (0.576,2.09);
      \node[right, font=\scriptsize] at (0.70,1.90) {600 a 1.600};
      \node[left, text width=3.5cm, align=right, font=\scriptsize]
        at (-0.15,1.10) {Carril exclusivo de autobús};
      \fill[cTP] (0,0.91) rectangle (2.88,1.29);
      \node[right, font=\scriptsize] at (3.00,1.10) {hasta 8.000};
      \node[left, text width=3.5cm, align=right, font=\scriptsize]
        at (-0.15,0.30) {Corredor segregado de transporte};
      \fill[cTP!70] (0,0.11) rectangle (9.00,0.49);
      \node[right, font=\scriptsize] at (9.12,0.30) {hasta 25.000};
    \end{scope}
  \end{tikzpicture}
  \caption{Las dos restricciones que ninguna mejora del vehículo levanta.
  Panel A: identidad geométrica, no medición; las velocidades son ilustrativas
  y el presupuesto de media hora por sentido procede de
  \textcite{zahavi1974} y \textcite{marchetti1994}. Panel B: capacidades por
  carril y sentido de la guía de diseño de calles con transporte público de
  \textcite{nacto2016}, que estima entre 600 y 1.600 personas por hora para un
  carril de tráfico general, hasta 8.000 para un carril exclusivo de autobús y
  hasta 25.000 para un corredor segregado. Elaboración propia.}
  \label{fig:geometria}
\end{figure}

Aquí es donde el debate suele desviarse. La respuesta habitual a los problemas
del automóvil es una mejora del automóvil: eléctrico, autónomo, conectado,
compartido. Es una respuesta razonable para una de las externalidades y
completamente inocua para las demás, y la Figura~\ref{fig:geometria} explica
por qué.

\textbf{La electrificación resuelve el tubo de escape y nada más.} En Costa
Rica el transporte genera el 41,55 por ciento de las emisiones nacionales de
gases de efecto invernadero y el 75,4 por ciento de las del sector energía, y
es la categoría que más ha crecido desde 1990, con un 243 por ciento
\parencite{dcc2021,pnud2022}. Cambiar la flota corrige ese problema, y por eso
el Plan Nacional de Descarbonización lo prioriza \parencite{descarbonizacion2019}.
Pero un vehículo eléctrico ocupa el mismo espacio en el carril, exige el mismo
aparcamiento y produce prácticamente la misma congestión. En 2026 el país
supera las \num{47000} unidades eléctricas \parencite{infobae2026}, en torno al
2 por ciento de una flota de unos dos millones: incluso al cien por cien, la
geometría del Panel B no se movería un milímetro.

\textbf{La geometría es implacable y no depende del motor.} Un carril dedicado
a tráfico general mueve entre 600 y 1.600 personas por hora; el mismo carril
dedicado a autobuses mueve hasta 8.000, y como corredor segregado hasta 25.000
\parencite{nacto2016}. La diferencia es de un factor de quince a cuarenta, y no
la produce ninguna innovación de producto: la produce la decisión de repartir
el espacio. A esto se añade que el activo está inmovilizado casi todo el
tiempo: el automóvil medio británico permanece aparcado el 96 por ciento de
las horas, el 80 en el domicilio y el 16 fuera de él \parencite{bates2012}.
Ninguna tecnología de propulsión mejora esa cifra de utilización.

\textbf{Y la demanda inducida cierra el bucle.} Si la mejora consiste en hacer
más agradable o más barato conducir, el resultado previsible según
\textcite{duranton2011} no es menos tráfico, sino más kilómetros recorridos.
El vehículo autónomo, tema del foro anterior de este mismo módulo, es el caso
extremo: al eliminar el coste de conducir, elimina también el principal freno
al kilometraje.

\section{Mi entorno: la Gran Área Metropolitana de Costa Rica}

\begin{figure}[ht]
  \centering
  \begin{tikzpicture}
    % ---------------- panel A: flota frente a autobús ----------------
    \node[font=\scriptsize\itshape, gray!35!black, right] at (-0.10,5.55)
      {Panel A. Dos series que se separan, base 100 en su primer dato};
    \foreach \yv/\lab in {0/0, 0.96/50, 1.92/100, 2.88/150, 3.84/200,
                          4.81/250} {
      \draw[gray!20] (0,\yv) -- (12.20,\yv);
      \node[left, font=\tiny, gray!55!black] at (-0.05,\yv) {\lab};
    }
    \draw[gray!60!black] (0,0) -- (12.20,0);
    \foreach \xv/\lab in {0/2005, 2.86/2010, 5.71/2015, 8.57/2020,
                          11.43/2025} {
      \draw[gray!60!black] (\xv,0) -- (\xv,-0.13);
      \node[below, font=\tiny, gray!55!black] at (\xv,-0.13) {\lab};
    }
    \draw[very thick, cAuto] (0,1.92) -- (11.43,4.81);
    \fill[cAuto] (0,1.92) circle (2pt);
    \fill[cAuto] (11.43,4.81) circle (2pt);
    \node[font=\tiny, cAuto, above left] at (0.35,1.98)
      {\num{751332} vehículos (2005)};
    \node[font=\tiny, cAuto, above left, align=right] at (11.35,4.90)
      {\num{1879790} vehículos (2025)\\[-0.25em] con derecho de circulación:
       \textbf{+150\,\%}};
    \draw[very thick, cTP, dashed] (4.00,1.92) -- (10.86,1.11);
    \fill[cTP] (4.00,1.92) circle (2pt);
    \fill[cTP] (10.86,1.11) circle (2pt);
    \node[font=\tiny, cTP, above right] at (4.10,1.95) {uso del autobús (2012)};
    \node[font=\tiny, cTP, below left, align=right] at (10.80,1.05)
      {\textbf{-42\,\%} en 2024};
    \node[font=\tiny, gray!45!black, align=left, right] at (0.15,0.55)
      {Trazo esquemático: sólo están documentados los extremos de cada serie,
       no su forma intermedia};
    % ---------------- panel B: calendario del tren eléctrico ----------
    \begin{scope}[shift={(0,-3.70)}]
      \node[font=\scriptsize\itshape, gray!35!black, right] at (-0.10,2.72)
        {Panel B. El calendario que decide cuándo puede cambiar mi respuesta};
      \draw[gray!60!black] (0,0.85) -- (12.20,0.85);
      \foreach \xv/\lab in {0/2025, 1.74/2026, 3.49/2027, 5.23/2028,
                            6.97/2029, 8.71/2030, 10.46/2031, 12.20/2032} {
        \draw[gray!60!black] (\xv,0.85) -- (\xv,0.72);
        \node[below, font=\tiny, gray!55!black] at (\xv,0.72) {\lab};
      }
      \fill[cTP!22] (5.23,0.90) rectangle (8.71,1.20);
      \node[font=\tiny, cTP!80!black] at (6.97,1.05) {diseño y construcción};
      \foreach \xv/\lab in {1.48/{créditos por USD 800\\millones al Congreso},
                            4.35/{adjudicación\\prevista}} {
        \draw[cAuto, thick] (\xv,0.85) -- (\xv,0.35);
        \fill[cAuto] (\xv,0.35) circle (1.7pt);
        \node[below, font=\tiny, cAuto, align=center] at (\xv,0.30) {\lab};
      }
      \draw[cTP, very thick] (2.44,0.85) -- (2.44,1.95);
      \fill[cTP] (2.44,1.95) circle (2.0pt);
      \node[above, font=\tiny, cTP, align=center] at (2.44,2.00)
        {ley firmada\\25 de mayo de 2026};
      \draw[cTP, thick] (2.90,0.85) -- (2.90,1.45);
      \fill[cTP] (2.90,1.45) circle (1.7pt);
      \node[right, font=\tiny, cTP, align=left] at (2.98,1.50)
        {licitación internacional};
      \foreach \xv/\lab in {8.71/{línea 2\\San José--Alajuela},
                            10.46/{línea 1\\San José--Paraíso}} {
        \draw[cPie, very thick] (\xv,0.85) -- (\xv,1.70);
        \fill[cPie] (\xv,1.70) circle (2.0pt);
        \node[above, font=\tiny, cPie, align=center] at (\xv,1.75) {\lab};
      }
      \draw[red!65!black, thick, densely dashed] (2.79,-0.35) -- (2.79,0.85);
      \node[below, font=\tiny, red!65!black, align=center] at (2.79,-0.37)
        {hoy};
      \draw[{Stealth[length=1.8mm]}-{Stealth[length=1.8mm]}, red!65!black,
        thick] (2.79,-1.00) -- (10.46,-1.00);
      \node[font=\tiny, red!65!black, above] at (6.60,-0.98)
        {cinco años hasta el sistema completo};
    \end{scope}
  \end{tikzpicture}
  \caption{Panel A: los vehículos con derecho de circulación pasaron de
  \num{751332} en 2005 a \num{1879790} en 2025, un crecimiento del 150 por
  ciento \parencite{observador2025}, mientras el uso del autobús cayó un 42 por
  ciento entre 2012 y 2024 \parencite{allen2026}; una medición independiente de
  la Universidad Nacional cifra la caída de pasajeros en 42,2 por ciento entre
  octubre de 2018 y mayo de 2025, con un 24,2 por ciento menos de operadores y
  cerca de cien rutas abandonadas \parencite{una2026}. Sólo se conocen los
  extremos de ambas series, de modo que el trazo intermedio es esquemático.
  Panel B: hitos del Tren Eléctrico de Pasajeros de la GAM, con 51,5 kilómetros
  de doble vía, 28 trenes y una capacidad proyectada superior a \num{100000}
  pasajeros diarios \parencite{semanario2025,incofer2026,teletica2026}.
  Elaboración propia.}
  \label{fig:costarica}
\end{figure}

Mi entorno está en el tramo alto de la escalera, y empeorando. La
Figura~\ref{fig:costarica} muestra las dos series que definen el problema:
la flota crece y el autobús se vacía al mismo tiempo. Es la señal más clara de
que no estamos ante una elección de consumo sino ante una degradación de la
alternativa, porque los dos movimientos se refuerzan: menos pasajeros implican
menos frecuencia, menos frecuencia implica menos pasajeros, y unas cien rutas
ya han desaparecido \parencite{una2026}.

El coste agregado está cuantificado y es de primer orden. El Programa Estado de
la Nación estimó que el tiempo perdido por congestión en la GAM equivale al 3,8
por ciento del producto interno bruto, del orden de \num{2500} millones de
dólares al año \parencite{semanario2018,allen2026}; la Universidad Nacional
eleva la cifra al 4 por ciento y la traduce a 400 millones de horas laborales
anuales \parencite{una2026}. En el plano doméstico, las familias
costarricenses destinan el 15,6 por ciento de su ingreso mensual a
transportarse \parencite{allen2024}. Y en el plano humano, 2025 cerró con 572
muertes en sitio por siniestros de tránsito, la cifra más alta de la serie
histórica, con los motociclistas concentrando el 46 por ciento de las víctimas
\parencite{observador2026}; la Organización Mundial de la Salud sitúa el total
mundial en 1,19 millones de muertes anuales y recuerda que son la principal
causa de muerte entre los 5 y los 29 años \parencite{who2023}.

Frente a eso, la respuesta institucional es la correcta en el papel: 51,5
kilómetros de tren eléctrico a doble vía entre Paraíso de Cartago y Alajuela,
28 trenes, 30 estaciones y una capacidad proyectada superior a \num{100000}
pasajeros diarios, financiados con unos 800 millones de dólares
\parencite{semanario2025}. Pero el Panel B de la figura contiene la parte
incómoda: la ley se firmó en mayo de 2026, la licitación empieza este año, la
adjudicación se espera en 2027 y la operación completa no llega antes de 2031
\parencite{teletica2026}. Cinco años es el plazo mínimo, y sólo si nada se
atrasa.

\section{¿Qué haría cambiar mi respuesta?}

Formulada como pregunta de gestión tecnológica, la cuestión del foro tiene una
respuesta operativa: mi respuesta cambia cuando se cumplan condiciones
medibles, no cuando aparezca un vehículo mejor. Enumero las que considero
necesarias, con el orden de magnitud que las haría creíbles.

\begin{enumerate}
  \item \textbf{Que el destino esté dentro del presupuesto de tiempo.} Si el
    trayecto puerta a puerta en transporte público supera en más del 50 por
    ciento al del automóvil, la sustitución no ocurre, por convicción que
    tenga el usuario. Es la consecuencia directa de
    \textcite{zahavi1974,marchetti1994}.
  \item \textbf{Que la densidad y la mezcla de usos alrededor de las
    estaciones cambien.} Un tren en una ciudad dispersa transporta poca gente,
    porque el problema no es el tramo central del viaje sino el primero y el
    último kilómetro. Las elasticidades de \textcite{ewing2010} y las tres
    \enquote{d} de \textcite{cervero1997} indican dónde hay que actuar:
    densidad, diversidad y diseño en el entorno de cada estación. Sin un plan
    regulador que lo permita, el tren será una obra cara con demanda pobre.
  \item \textbf{Que el espacio se reasigne, no sólo se añada.} Añadir capacidad
    sin retirar espacio al automóvil no reduce la congestión
    \parencite{duranton2011}. El factor de quince a cuarenta del Panel B de la
    Figura~\ref{fig:geometria} sólo se materializa si los carriles se dedican.
  \item \textbf{Que la alternativa deje de degradarse.} Mientras el autobús
    siga perdiendo pasajeros al ritmo del 42 por ciento por década, cada mejora
    del tren llegará a un sistema con menos capilaridad de la que tenía cuando
    se diseñó.
\end{enumerate}

El horizonte razonable, si todo lo anterior avanza, es el de la ciudad de
proximidad que describe \textcite{moreno2021}: no un mundo sin automóviles,
sino un mundo en el que el automóvil deja de ser la condición de acceso al
empleo. Nótese que esa es exactamente la situación que ya viven el 37,6 por
ciento de los hogares de la almendra de Madrid y el 56,7 por ciento de los
hogares de Nueva York \parencite{crtm2019,acs2024}.

\section{Precisiones y limitaciones}

\begin{enumerate}
  \item \textbf{Las dos métricas de la Figura~\ref{fig:escalera} no son
    intercambiables.} La cuota sobre todos los viajes y la cuota sobre los
    desplazamientos al trabajo miden cosas distintas: la segunda excluye los
    viajes cortos de compra y acompañamiento, que son los más fáciles de hacer
    a pie, y por eso tiende a dar valores más altos para el automóvil. He
    conservado ambas porque las fuentes disponibles para las cuatro ciudades
    del enunciado no publican la misma métrica, y las he distinguido con color
    para que la comparación no se lea como si fuera homogénea. El ordenamiento
    general, sin embargo, es robusto: los extremos están separados por un
    factor de 3,5, muy por encima del sesgo que introduce la diferencia de
    definición.
  \item \textbf{El censo australiano de 2021 se levantó en pandemia.} El 10 de
    agosto de 2021 varias capitales australianas estaban confinadas, lo que
    infló el trabajo desde casa hasta el 21 por ciento nacional
    \parencite{abs2021}. Australia del Sur no estaba confinada ese día, de modo
    que la cifra de Adelaida es de las menos afectadas, pero conviene tomarla
    como orden de magnitud.
  \item \textbf{Los datos de Madrid son de 2018 y los de Nueva York y Los
    Ángeles de 2024.} Entre ambas fechas está la pandemia, que en Estados
    Unidos elevó el trabajo desde casa al 12,5 por ciento en Nueva York y al
    16,0 en Los Ángeles \parencite{acs2024}. Eso reduce mecánicamente todas las
    cuotas de desplazamiento, incluida la del automóvil, y por tanto juega en
    contra de mi argumento, no a favor: la diferencia real entre las dos
    ciudades estadounidenses es probablemente mayor que la medida.
  \item \textbf{La flota costarricense se cuenta de tres maneras distintas.}
    El Instituto Nacional de Seguros puso al cobro \num{1879790} derechos de
    circulación para 2025 \parencite{observador2025}, el registro total de
    automotores supera los 3,1 millones porque no se desinscriben las unidades
    en desuso, y \textcite{allen2026} trabaja con unos 2,1 millones de
    vehículos registrados. He usado la serie de derechos de circulación por ser
    la única que se aproxima a la flota efectivamente en uso, y advierto que la
    comparación entre fuentes no es directa.
  \item \textbf{La afirmación de que San José tiene el segundo peor tránsito
    del mundo no es utilizable.} Circula ampliamente, pero procede del índice
    de Numbeo, que se construye con respuestas voluntarias de usuarios y no con
    mediciones de sondas de tráfico; el índice de TomTom, que sí las usa, no
    cubre San José por falta de datos granulares. Lo que sí resiste el
    escrutinio son las estimaciones de coste del Programa Estado de la Nación y
    de la Universidad Nacional, que proceden de encuestas de empleo y de
    matrices de origen y destino.
  \item \textbf{El Panel A de la Figura~\ref{fig:geometria} es una identidad
    geométrica, no una medición.} Las velocidades puerta a puerta reales son
    menores que las nominales y el presupuesto de tiempo de viaje es estable,
    no rigurosamente constante: el propio \textcite{zahavi1974} advirtió que
    debe medirse en cada lugar antes de usarse. El argumento no depende del
    valor exacto, sino de que el tiempo esté acotado y la velocidad no.
  \item \textbf{No puedo contrastar el contrafactual.} La tesis de que la forma
    urbana determina la dependencia del automóvil es coherente con la teoría y
    con todas las mediciones que he podido reunir, pero la causalidad podría
    ir, al menos en parte, en sentido contrario: es posible que quien prefiere
    no conducir se mude al centro. \textcite{ewing2010} discuten esa
    autoselección residencial y concluyen que atenúa las elasticidades pero no
    las anula. Con datos agregados como los míos, el efecto no es separable.
\end{enumerate}

\section*{Conclusiones}

\textbf{¿Puedo imaginarme una vida sin vehículo privado?} Puedo imaginarla con
precisión, porque está documentada: la viven el 37,6 por ciento de los hogares
del centro de Madrid \parencite{crtm2019}, el 56,7 por ciento de los de Nueva
York \parencite{acs2024} y, en la práctica
cotidiana, el 49,8 por ciento de los viajes que en mi propio país se hacen en
autobús. Lo que no puedo es vivirla yo hoy, en la GAM, sin renunciar a una
parte sustancial de mi acceso al trabajo y a los servicios.

\textbf{Y esa asimetría es el punto.} La diferencia entre unos y otros no está
en el vehículo, que es el mismo en todas partes, ni en la convicción ambiental,
que se distribuye de forma bastante pareja. Está en cuántos destinos caben
dentro de media hora sin coche, y eso lo fijó quien decidió el uso del suelo
hace veinte, cuarenta o setenta años. La escalera del 20 al 71 por ciento de la
Figura~\ref{fig:escalera} no mide preferencias: mide urbanismo.

La lección que me llevo como CTO se puede formular en una frase:
\textbf{cuando el desempeño observado de una tecnología varía por un factor de
tres según el entorno donde se despliega, el objeto de la gestión ya no es la
tecnología, es el entorno}. Es el mismo error de nivel que discutimos en el
foro del módulo anterior a propósito de IPv6: optimizar la figura de mérito del
artefacto mientras la que decide el resultado es la del sistema. Aquí el sistema se llama
densidad, mezcla de usos y reparto del espacio, y su constante de tiempo se
mide en décadas, no en ciclos de producto.

\textbf{Pregunta para el foro.} Me interesa contrastar el dato con el resto del
grupo, porque venimos de entornos muy distintos: en su ciudad, ¿cuánto tarda
puerta a puerta el mismo trayecto en transporte público y en automóvil? Si la
razón entre ambos supera el 1,5, mi hipótesis es que ustedes también conducen,
con independencia de lo que opinen del automóvil. Y una segunda pregunta para
quienes trabajan en industrias reguladas: ¿han visto algún caso en el que la
mejora del producto haya conseguido, por sí sola, cambiar un resultado que en
realidad estaba determinado por la infraestructura?

\printbibliography[title={Referencias bibliográficas}]

\end{document}
